\begin{tikzpicture}[
    font=\sffamily\scriptsize,
    >={Latex[length=2.8mm]},
    flow/.style={->,very thick,draw=black!70},
    title/.style={font=\sffamily\small\bfseries,align=center},
    note/.style={align=center,text=black!75},
    index/.style={font=\sffamily\normalsize,text=black!85,inner sep=0pt}
]
% Spatial tokens and a fixed two-window partition.
\node[title] at (2.0,2.85) {cuadrícula de tokens};
\foreach \row in {0,1}{
  \foreach \col in {0,1,2,3}{
    \pgfmathtruncatemacro{\token}{4*\row+\col+1}
    \ifnum\col<2
      \def\cellcolor{blue!18}
    \else
      \def\cellcolor{orange!20}
    \fi
    \draw[fill=\cellcolor,draw=black!55,line width=0.7pt]
      (\col,1-\row) rectangle ++(1,1);
    \node at (\col+0.5,1.5-\row) {$x_{\token}$};
  }
}
\draw[red!75!black,line width=2.1pt] (0,1) rectangle ++(1,1);
\node[note] at (2.0,-0.48) {dos ventanas de cuatro tokens;\\$x_1$ es la query señalada};

\draw[flow] (4.45,2.55) -- (5.25,2.55);

% Dense attention matrix.
\node[title] at (7.05,2.85) {atención global};
\foreach \i in {0,...,7}{
  \foreach \j in {0,...,7}{
    \pgfmathtruncatemacro{\shade}{12+mod(17*\i+23*\j,72)}
    \draw[fill=blue!\shade,draw=white,line width=0.45pt]
      (5.55+0.38*\j,2.22-0.38*\i) rectangle ++(0.38,-0.38);
  }
}
\draw[red!75!black,line width=1.4pt] (5.55,2.22) rectangle ++(3.04,-0.38);
% The same permutation is used on both matrix axes and in both panels.
\foreach \position/\token in {0/1,1/2,2/5,3/6,4/3,5/4,6/7,7/8}{
  \node[index,anchor=south] at (5.74+0.38*\position,2.35) {\token};
  \node[index,anchor=east] at (5.42,2.03-0.38*\position) {\token};
}
\node[rotate=90,anchor=south] at (4.98,0.70) {queries};
\node at (7.07,-1.06) {keys};
\node[note] at (7.07,-1.55) {$8\times8$: cada token compara\\con todos los demás};

\draw[flow] (9.05,2.55) -- (9.85,2.55);

% Block-diagonal window attention map.
\node[title] at (11.68,2.85) {atención por ventanas};
\foreach \i in {0,...,7}{
  \foreach \j in {0,...,7}{
    \pgfmathtruncatemacro{\blocki}{int(\i/4)}
    \pgfmathtruncatemacro{\blockj}{int(\j/4)}
    \ifnum\blocki=\blockj
      \pgfmathtruncatemacro{\shade}{18+mod(19*\i+29*\j,70)}
      \def\cellcolor{blue!\shade}
    \else
      \def\cellcolor{black!8}
    \fi
    \draw[fill=\cellcolor,draw=white,line width=0.45pt]
      (10.16+0.38*\j,2.22-0.38*\i) rectangle ++(0.38,-0.38);
  }
}
\draw[red!75!black,line width=1.4pt] (10.16,2.22) rectangle ++(3.04,-0.38);
\draw[black!65,line width=1pt] (10.16,2.22) rectangle ++(1.52,-1.52);
\draw[black!65,line width=1pt] (11.68,0.70) rectangle ++(1.52,-1.52);
\foreach \position/\token in {0/1,1/2,2/5,3/6,4/3,5/4,6/7,7/8}{
  \node[index,anchor=south] at (10.35+0.38*\position,2.35) {\token};
  \node[index,anchor=east] at (10.03,2.03-0.38*\position) {\token};
}
\node[rotate=90,anchor=south] at (9.59,0.70) {queries};
\node at (11.68,-1.06) {keys};
\node[note] at (11.68,-1.55) {bloques activos: pesos locales;\\gris: conexiones enmascaradas};

% Weight legend.
\foreach \k in {0,...,7}{
  \pgfmathtruncatemacro{\shade}{8+11*\k}
  \fill[blue!\shade] (13.75+0.24*\k,1.34) rectangle ++(0.24,0.28);
}
\draw[black!45] (13.75,1.34) rectangle ++(1.92,0.28);
\node[anchor=north west] at (13.70,1.22) {peso bajo};
\node[anchor=north east] at (15.72,1.22) {alto};
\node[note,text width=2.2cm] at (14.72,0.10)
    {cada fila se normaliza por separado};
\end{tikzpicture}
